% Source: Weinberg GR, physical PDF pp. 121--123 (printed pp. 96--98).
% Coverage: Section 4.3, Tensor Algebra; no numbered equations.
% Figures/tables/footnotes: none.
% Status: source-reviewed and compile-clean.
% Uncertainties: none.

\subsection{Tensor Algebra}\label{sec:4.3}

The next step in our program of constructing equations invariant under general
coordinate transformations is to learn how to put tensors together to form
other tensors. This is accomplished through a few simple algebraic operations.

\medskip
\noindent\textbf{(A) Linear Combinations.}
A linear combination of tensors with the same upper and lower indices is a
tensor with these indices. For instance, let \(A^\mu{}_\nu\) and
\(B^\mu{}_\nu\) be mixed tensors, and let
\[
  T^\mu{}_\nu\equiv aA^\mu{}_\nu+bB^\mu{}_\nu,
\]
where \(a\) and \(b\) are scalars. Then \(T^\mu{}_\nu\) is a tensor, because
\begin{align*}
  T'^\mu{}_\nu
  &=aA'^\mu{}_\nu+bB'^\mu{}_\nu
  \\
  &=
  a\frac{\partial x'^\mu}{\partial x^\rho}
   \frac{\partial x^\sigma}{\partial x'^\nu}A^\rho{}_\sigma
  +b\frac{\partial x'^\mu}{\partial x^\rho}
   \frac{\partial x^\sigma}{\partial x'^\nu}B^\rho{}_\sigma
  \\
  &=
  \frac{\partial x'^\mu}{\partial x^\rho}
  \frac{\partial x^\sigma}{\partial x'^\nu}T^\rho{}_\sigma.
\end{align*}

\medskip
\noindent\textbf{(B) Direct Products.}
The product of the components of two tensors yields a tensor whose upper and
lower indices consist of all the upper and lower indices of the two original
tensors. For instance, if \(A^\mu{}_\nu\) and \(B^\rho\) are tensors, and
\[
  T^{\mu\rho}{}_\nu\equiv A^\mu{}_\nu B^\rho,
\]
then \(T^{\mu\rho}{}_\nu\) is a tensor; that is,
\begin{align*}
  T'^{\mu\rho}{}_\nu
  &=A'^\mu{}_\nu B'^\rho
  \\
  &=
  \frac{\partial x'^\mu}{\partial x^\lambda}
  \frac{\partial x^\kappa}{\partial x'^\nu}A^\lambda{}_\kappa
  \frac{\partial x'^\rho}{\partial x^\sigma}B^\sigma
  \\
  &=
  \frac{\partial x'^\mu}{\partial x^\lambda}
  \frac{\partial x^\kappa}{\partial x'^\nu}
  \frac{\partial x'^\rho}{\partial x^\sigma}
  T^{\lambda\sigma}{}_\kappa.
\end{align*}

\medskip
\noindent\textbf{(C) Contraction.}
Setting an upper and lower index equal and summing it over its four values
yields a new tensor with these two indices absent. For instance, if
\(T^\mu{}_\nu{}^{\rho\sigma}\) is a tensor and
\[
  T^{\mu\rho}\equiv T^\mu{}_\nu{}^{\rho\nu},
\]
then \(T^{\mu\rho}\) is a tensor; that is,
\begin{align*}
  T'^{\mu\rho}
  &=T'^\mu{}_\nu{}^{\rho\nu}
  \\
  &=
  \frac{\partial x'^\mu}{\partial x^\kappa}
  \frac{\partial x^\lambda}{\partial x'^\nu}
  \frac{\partial x'^\rho}{\partial x^\eta}
  \frac{\partial x'^\nu}{\partial x^\tau}
  T^\kappa{}_\lambda{}^{\eta\tau}
  \\
  &=
  \frac{\partial x'^\mu}{\partial x^\kappa}
  \frac{\partial x'^\rho}{\partial x^\eta}
  T^\kappa{}_\lambda{}^{\eta\lambda}
  \\
  &=
  \frac{\partial x'^\mu}{\partial x^\kappa}
  \frac{\partial x'^\rho}{\partial x^\eta}T^{\kappa\eta}.
\end{align*}

These three operations can of course be combined in various ways. One
particularly important combined operation results in the \textit{raising and
lowering of indices}. If we take the direct product of a contravariant or
mixed tensor \(T\) with the metric tensor \(g_{\mu\nu}\), and contract the
index \(\mu\) with one of the contravariant indices of \(T\), we get a new
tensor in which this contravariant index is replaced by a covariant index
\(\nu\). For instance, if \(T^{\mu\rho}{}_\sigma\) is a tensor, and we define
\[
  S_\nu{}^\rho{}_\sigma\equiv g_{\mu\nu}T^{\mu\rho}{}_\sigma,
\]
then by rules (B) and (C), \(S_\nu{}^\rho{}_\sigma\) will be a tensor.
Similarly, if we take the direct product of a covariant or mixed tensor \(T\)
with the inverse metric tensor \(g^{\mu\nu}\), and contract the index \(\mu\)
with one of the covariant indices of \(T\), we get a new tensor in which this
covariant index is replaced by a contravariant index \(\nu\). For instance,
if \(S_\mu{}^\rho{}_\sigma\) is a tensor, and we define
\[
  R^{\nu\rho}{}_\sigma\equiv g^{\mu\nu}S_\mu{}^\rho{}_\sigma,
\]
then \(R^{\nu\rho}{}_\sigma\) is also a tensor. Note that lowering an index
and then raising it again gives back the original tensor; for instance, in
the examples cited above, we lowered an index on \(T\) to get \(S\) and then
raised it again to get \(R\), so \(R=T\), because
\begin{align*}
  R^{\nu\rho}{}_\sigma
  &\equiv g^{\mu\nu}S_\mu{}^\rho{}_\sigma
   \equiv g^{\mu\nu}g_{\mu\lambda}T^{\lambda\rho}{}_\sigma
  \\
  &=\delta^\nu{}_\lambda T^{\lambda\rho}{}_\sigma
   =T^{\nu\rho}{}_\sigma.
\end{align*}

By raising and lowering indices we can write a tensor with \(N\) indices in
\(2^N\) different ways. Since these are all physically equivalent, it is
customary to use the same symbol for all \(2^N\) tensors, distinguishing them
only by their index locations.

For the sake of completeness, it should be mentioned that the tensor obtained
by raising one index on the metric tensor \(g_{\mu\nu}\), or by lowering one
index on the inverse metric tensor \(g^{\mu\nu}\), is precisely the Kronecker
tensor, because
\[
  g^{\mu\lambda}g_{\lambda\nu}=\delta^\mu{}_\nu.
\]
Also, raising both indices on \(g_{\mu\nu}\) gives the inverse tensor,
\[
  g^{\lambda\mu}g^{\sigma\nu}g_{\mu\nu}
  =g^{\lambda\mu}\delta^\sigma{}_\mu
  =g^{\lambda\sigma},
\]
and lowering both indices on \(g^{\mu\nu}\) gives the metric tensor
\(g_{\mu\nu}\).

The reader will have noticed that this discussion of tensor algebra is
precisely the same as the corresponding discussion in the chapter on special
relativity (see Section~\ref{sec:2.5}), with one important exception. I have
not yet mentioned differentiation. This is because the derivative of a tensor
is in general not a tensor. In Section~\ref{sec:4.6} we shall see that there
is a kind of differentiation, called covariant differentiation, that provides
one more way of constructing tensors from other tensors.
